% !TeX program = xelatex
\documentclass[14pt,a4paper]{extarticle}
\usepackage{fontspec}
\setmainfont{Times New Roman}
\setmonofont{DejaVuSansMono.ttf}[BoldFont=DejaVuSansMono-Bold.ttf,Scale=0.82]
\usepackage[english]{babel}
\babelprovide[import,main]{russian}
\usepackage{amsmath,unicode-math}
\setmathfont{texgyretermes-math.otf}
\usepackage[left=30mm,right=15mm,top=20mm,bottom=20mm]{geometry}
\usepackage{setspace,indentfirst}
\onehalfspacing
\setlength{\parindent}{1.25cm}
\setlength{\parskip}{0pt}
\usepackage{titlesec}
\titleformat{\section}{\normalfont\bfseries}{\thesection}{0.6em}{}
\titleformat{\subsection}{\normalfont\bfseries}{\thesubsection}{0.6em}{}
\titlespacing*{\section}{\parindent}{0pt}{12pt}
\titlespacing*{\subsection}{\parindent}{12pt}{6pt}
\newcommand{\sectionbreak}{\clearpage}
\usepackage{graphicx,float,placeins,caption,array,tabularx}
\DeclareCaptionLabelSeparator{dash}{~---~}
\captionsetup{labelsep=dash,font=normalsize,skip=6pt}
\captionsetup[figure]{name=Рисунок,justification=centering}
\captionsetup[table]{name=Таблица,justification=raggedright,singlelinecheck=false,position=top}
\renewcommand{\arraystretch}{1.15}
\newcolumntype{L}{>{\raggedright\arraybackslash}X}
\usepackage{fvextra,newfloat,needspace}
\DeclareFloatingEnvironment[name=Листинг]{listing}
\captionsetup[listing]{hypcap=false,justification=raggedright,singlelinecheck=false}
\fvset{numbers=left,numbersep=5pt,fontsize=\scriptsize,frame=single,
  breaklines=true,breakanywhere=true,breaksymbolleft={},tabsize=4}
\usepackage[hidelinks]{hyperref}
\newcommand{\code}[1]{\texttt{\detokenize{#1}}}
\newsavebox{\functiondescription}
\newcommand{\fn}[3]{\par
  \setbox\functiondescription=\vbox{\hsize=\linewidth
  \noindent\textbf{\code{#1}}. \textit{Вход:} #2\par
  \noindent\textit{Выход:} #3\par}
  \Needspace{\dimexpr\ht\functiondescription+\dp\functiondescription+6\baselineskip\relax}
  \medskip\unvbox\functiondescription}
\newcommand{\structsection}[1]{\clearpage\phantomsection
  \addcontentsline{toc}{section}{#1}
  \begin{center}\textbf{#1}\end{center}\vspace{6pt}}
\newcommand{\consolefig}[3]{\begin{figure}[H]\centering
  \includegraphics[width=#2\textwidth]{fig/#1.png}
  \caption{#3}\label{fig:#1}\end{figure}}
\newcommand{\functioncode}[2]{\par\Needspace{5\baselineskip}\smallskip
  \begin{singlespace}\captionof{listing}{\code{#1}}
  \VerbatimInput{functions/#2}\end{singlespace}}
\newcommand{\aA}{a_{Ai}(x_i)}
\newcommand{\aB}{a_{Bi}(x_i)}
\newcommand{\aU}{a_{\hat{U}}(x_i)}
\pagestyle{plain}
\sloppy

\begin{document}
\begin{titlepage}
\begin{singlespace}
\begin{center}
Министерство науки и высшего образования Российской Федерации\\
Санкт-Петербургский политехнический университет Петра Великого\\
Институт компьютерных наук и кибербезопасности\\
Высшая школа технологий искусственного интеллекта
\end{center}
\vspace{5cm}
\begin{center}
\textbf{ОТЧЁТ}\\
по лабораторной работе № 1\\[6pt]
«Генерация бинарного кода Грея. Мультимножества и операции над ними»\\[6pt]
по дисциплине «Дискретная математика»
\end{center}
\vspace{3.5cm}
\begin{flushright}
\begin{tabular}{@{}l@{}}
Выполнил\\
студент гр. 5130201/50302 \underline{\hspace{2.5cm}} Р. Ю. Дужнов\\[12pt]
Проверил\\
преподаватель \underline{\hspace{2.5cm}} А. В. Востров\\[12pt]
« \underline{\hspace{1cm}} » \underline{\hspace{3cm}} 2026 г.
\end{tabular}
\end{flushright}
\vfill
\begin{center}Санкт-Петербург\\2026\end{center}
\end{singlespace}
\end{titlepage}
\setcounter{page}{2}
\tableofcontents

\structsection{ВВЕДЕНИЕ}

В работе требуется разработать программу для работы с мультимножествами, элементами которых служат слова бинарного кода Грея.

Программа должна генерировать бинарный код Грея разрядности $n$, которую задаёт пользователь, и заполнять универсум мультимножеств его словами. На основе универсума программа формирует два мультимножества $A$ и $B$, мощности которых задаёт пользователь, а способ заполнения каждого из них — вручную или автоматически — пользователь выбирает сам.

Над мультимножествами $A$ и $B$ программа должна вычислять и выводить результаты операций объединения, пересечения, разности, симметрической разности, дополнения каждого из мультимножеств до универсума, а так же арифметических суммы, разности, произведения и деления. Так же программа должна защищать работу от некорректного ввода и предоставлять пользователю меню.

\section{Математическое описание}
\begingroup
\setlength{\abovedisplayskip}{6pt plus 2pt minus 1pt}
\setlength{\belowdisplayskip}{6pt plus 2pt minus 1pt}
\subsection{Множества и мультимножества}

Множество - это любая определенная совокупность объектов. Объекты, из которых составлено множество, называются его элементами. Элементы множества различны и отличимы друг от друга.

Порождающее множество (универсум) - множество $U$, из элементов которого образуются все мультимножества этого семейства. Универсум состоит из $N=2^n$ слов бинарного кода Грея разрядности n. Наполнение универсума формируется с помощью порождающей процедуры [функция кода грея] на основе выбранной пользователем разрядности $n$, $n\geq 0$ и $n\leq 19$.

Мультимножеством $\hat{X}$ называется совокупность элементов множества $X$, в которую элемент $x_i$ входит $a_i$ раз, $a_i \geq  0$. Мультимножество обозначается следующим образом:
\[
\langle a_1(x_1), \dots, a_n(x_n) \rangle.
\]
Число $a_i$ называется показателем элемента $x_i$, множество $X$ - носителем, число $m=a_1 + ... + a_n$ - мощностью мультимножества.

У мультимножества $\hat{U}$, носителем которого является универсум, все элементы имеют одинаковые показатели $k$. Формирование показателей возможно двумя способами - либо случайным числом, либо заданным пользователем числом. При этом у всех элементов одинаковые показатели. $k\geq 1$, $k\leq 100$. Мощность мультимножества $\hat{U}$ равна $m_{\hat{U}}=k*2^n$

Элемент $x_i$ принадлежит мультимножеству $\hat{X}$, если его показатель $a_i$ строго больше нуля ($a_i>0$). При этом показатель любого элемента не превышает показателя соответствующего элемента в $\hat{U}$. В общем случае, мультимножество B включает мультимножество A, $A \subseteq B$, если $a_{Ai}(x_i)\leq a_{Bi}(x_i)$.

Заполнение мультимножеств $A$ и $B$ на основе $\hat{U}$ реализовано в двух вариантах. Первым вариантом является заполнение автоматически [рандом функция], вторым вариантом - ручное заполнение.

В данной работе мультимножество задается с помощью пары: универсума и вектора показателей, элементы которого строго соответствуют элементам универсума по порядковому номеру.

\subsection{Операции над мультимножествами}

В лабораторной работе теоретико-множественные и арифметические операции над мультимножествами заданы следующим образом.

Теоретико-множественные операции:

Объединение ($A \cup B$):
\[
A \cup B = \{a_i(x_i) \mid a_i(x_i) = \max\{a_{Ai}(x_i), a_{Bi}(x_i)\}\}
\]
Пересечение ($A \cap B$):
\[
A \cap B = \{a_i(x_i) \mid a_i(x_i) = \min\{a_{Ai}(x_i), a_{Bi}(x_i)\}\}
\]
Дополнение ($\overline{A}$):
\[
\overline{A} = \{a_i(x_i) \mid a_i(x_i) = a_{\hat{U}}(x_i) - a_{Ai}(x_i)\}
\]
Дополнение ($\overline{B}$):
\[
\overline{B} = \{a_i(x_i) \mid a_i(x_i) = a_{\hat{U}}(x_i) - a_{Bi}(x_i)\}
\]
Разность ($A \setminus B$):
\[
A \setminus B = \{a_i(x_i) \mid a_i(x_i) = \min\{a_{Ai}(x_i), a_{\hat{U}}(x_i) - a_{Bi}(x_i)\}\}
\]
Разность ($B \setminus A$):
\[
B \setminus A = \{a_i(x_i) \mid a_i(x_i) = \min\{a_{Bi}(x_i), a_{\hat{U}}(x_i) - a_{Ai}(x_i)\}\}
\]
Симметрическая разность ($A \triangle B$=$(A \setminus B) \cup (B \setminus A)$):
\[
\begin{aligned}
A \triangle B = \{a_i(x_i) \mid a_i(x_i) = \max\{&
\min\{a_{Ai}(x_i), a_{\hat{U}}(x_i) - a_{Bi}(x_i)\},\\
&\min\{a_{Bi}(x_i), a_{\hat{U}}(x_i) - a_{Ai}(x_i)\}\}\}
\end{aligned}
\]
Арифметические операции:

Арифметическая сумма ($A + B$):
\[
A + B = \{a_i(x_i) \mid a_i(x_i) = \min\{a_{Ai}(x_i) + a_{Bi}(x_i), a_{\hat{U}}(x_i)\}\}
\]
Арифметическая разность ($A - B$):
\[
A - B = \{a_i(x_i) \mid a_i(x_i) = \max\{a_{Ai}(x_i) - a_{Bi}(x_i), 0\}\}
\]
Арифметическое произведение ($A * B$):
\[
A * B = \{a_i(x_i) \mid a_i(x_i) = \min\{a_{Ai}(x_i) * a_{Bi}(x_i), a_{\hat{U}}(x_i)\}\}
\]
Арифметическое деление ($A : B$):
\[
A : B = \begin{cases}
\{a_i(x_i) \mid a_i(x_i) = \lfloor a_{Ai}(x_i) : a_{Bi}(x_i) \rfloor\}, & a_{Bi}(x_i) \neq 0 \\
0, & a_{Bi}(x_i) = 0
\end{cases}
\]
Результатом указанных операций над мультимножествами является мультимножество.

\subsection{Бинарный код Грея}

Бинарным кодом Грея разрядности n называется последовательность всех $2^n$ двоичных кодов длины n, в которой любые два соседних кода различаются ровно в одном разряде.

В данной работе алгоритм генерации бинарного кода Грея реализован следующим образом.

Алгоритм начинает с нуля и на каждой итерации $i$ ($i>0, i<m_{\hat{U}}$), инвертирует один бит. При $i=0$, код, состоящий из n нулей, сразу добавляется в последовательность. Чтобы определить, какой конкретно бит инвертировать, алгоритм определяет количество младших нулевых битов числа $i$, что соответствует позиции младшего единичного бита. Для этого выполняется проверка числа на четность ($i\%2=0$): пока оно четно, число делится на два $i/2$ и количество подсчитанных нулей увеличивается на единицу. Число нулей в конце двоичной записи $i$ - номер бита справа, считая с нуля, который требуется инвертировать.

В результате формируется универсум, состоящий из $2^n$ элементов.

\subsection{Алгоритм автоматического заполнения}

Алгоритм работает следующим образом. Создается пустое мультимножество с нулевыми показателями. Затем, $c$ раз ($c\leq m_{\hat{U}}$, $c\geq  1$, $c\leq k*2^n$) добавляется единица к показателю случайно выбранного элемента. При этом показатель элемента в формируемом мультимножестве не превышает показатель этого же элемента в $\hat{U}$. Если этот показатель достигается, в следующих итерациях элемент игнорируется.

\endgroup
\section{Особенности реализации}

\input{implementation.tex}

\section{Результаты работы}

Программа обрабатывает некорректный ввод и вычисляет операции по заданным формулам. В консоли \code{U} обозначает универсум слов, \code{Û} — мультимножество $\hat U$, \code{a_Ui}, \code{a_Ai}, \code{a_Bi} — показатели соответствующего слова, \code{m_X} — мощность мультимножества. Запись \code{⟨2(000), 1(011)⟩} означает два вхождения слова 000 и одно вхождение слова 011; нулевые показатели пропускаются. Обозначения \code{¬A}, \code{¬B} соответствуют дополнениям, \code{A * B} — произведению, \code{A : B} — арифметическому делению.

Проверен основной сценарий при $n=3$, $k=3$: $A$ мощности 10 заполнено автоматически, $B$ мощности 8 — вручную. Последовательность действий приведена на рисунках~\ref{fig:01_help}–\ref{fig:06_matrix}.

\consolefig{01_help}{1}{Список команд и справка}
\consolefig{02_universe}{1}{Создание универсума при $n = 3$, $k = 3$}
\consolefig{03_fill}{1}{Автоматическое заполнение $A$ ($m_A = 10$) и ручное заполнение $B$ ($m_B = 8$)}
\consolefig{04_table}{1}{Таблица показателей}
\consolefig{05_operations}{1}{Результаты операций и проверки}
\consolefig{06_matrix}{1}{Матрица показателей и выход}

Для ручной проверки выбраны слова 000, 001, 011. Их показатели в $A$ равны $(1,3,1)$, в $B$ — $(2,0,3)$, в $\hat U$ — $(3,3,3)$. Значения в таблице~\ref{tab:manual} совпали с соответствующими столбцами матрицы на рисунке~\ref{fig:06_matrix}.

\begin{table}[H]
\caption{Ручная проверка показателей результатов}
\label{tab:manual}
\centering
\begin{singlespace}
\begin{tabular}{|l|c|c|c|}
\hline
Операция & 000 & 001 & 011 \\\hline
$A\cup B$ & 2 & 3 & 3 \\\hline
$A\cap B$ & 1 & 0 & 1 \\\hline
$A\setminus B$ & 1 & 3 & 0 \\\hline
$B\setminus A$ & 2 & 0 & 2 \\\hline
$A\triangle B$ & 2 & 3 & 2 \\\hline
$\overline A$ & 2 & 0 & 2 \\\hline
$\overline B$ & 1 & 3 & 0 \\\hline
$A+B$ & 3 & 3 & 3 \\\hline
$A-B$ & 0 & 3 & 0 \\\hline
$A*B$ & 2 & 0 & 3 \\\hline
$A:B$ & 0 & 0 & 0 \\\hline
\end{tabular}
\end{singlespace}
\end{table}

Проверки мощностей также выполнены: $m_{\overline A}=24-10=14$, $m_{A\cup B}+m_{A\cap B}=16+2=10+8$. Например, для слова 011 показатель в $A\setminus B$ равен $\min\{1,3-3\}=0$, а в $A:B$ — $\lfloor1:3\rfloor=0$.

\structsection{ЗАКЛЮЧЕНИЕ}

В результате работы реализована программа для работы с мультимножествами, элементами которых служат слова бинарного кода Грея. Программа генерирует бинарный код Грея соответствующей разрядности, заполняем его словами универсум, на основе универсума формирует два мультимножества, вычисляет и выводит результаты операций объединения, пересечения, разности, симметрической разности, дополнения каждого из мультимножеств до универсума, а так же арифметических суммы, разности, произведения и деления. Так же программа обрабатывает некорректный ввод и предоставляет пользователю меню.

Достоинством программы является корректная обработка пограничных случаев: нулевой разрядности кода, пустых мультимножеств, предельных значений показателей и мощности, полного заполнения мультимножества и нулевого показателя делителя. Предусмотрены проверка пользовательского ввода и отмена заполнения с сохранением прежних данных.

Главным недостатком программы является большой расход памяти. В программе для каждого мультимножества хранятся показатели всех слов, даже если большинство из них нули. При увеличении разрядности кода на единицу, требуется в два раза больше памяти.

Программу можно масштабировать следующими функциями: проверка равенства через два включения, получение обычного множества элементов из мультимножества, умножение мультимножества на скаляр, декартово произведение. Для этих функций ядро программы изменять не потребуется.

\structsection{СПИСОК ИСПОЛЬЗОВАННЫХ ИСТОЧНИКОВ}
\begin{enumerate}
\item Новиков, Ф.\,А. Дискретная математика для программистов : учебник для вузов / Ф.\,А.~Новиков.~--- 3-е изд.~--- Санкт-Петербург : Питер, 2009.~--- 384~с.
\item Востров, А.\,В. Дискретная математика : курс лекций / А.\,В.~Востров; Санкт-Петербургский политехнический университет Петра Великого.~--- Санкт-Петербург, 2026.
\end{enumerate}

\end{document}
